\begin{abstract}
When one rotor of a quadrotor (a drone with four propellers) stops completely, the drone can no longer hover level. The only steady way it can still fly is to spin continuously about a slightly tilted axis. Mueller and D'Andrea called this state the \emph{relaxed hover} and showed that it can be controlled. This project studies, in simulation, how to bring a drone into this spinning state reliably after a rotor fails.

We ask two connected questions. \textbf{First}, the physics model of the spinning drone is good but not perfect: can we keep it and learn only its error from data, with the tools of Koopman methods, and does this predict and control better than the physics model alone or a Koopman model of the full dynamics? \textbf{Second}, a fault detector needs time to notice the failure: how long may this take before a safe recovery is no longer possible, and can such a limit be guaranteed mathematically?

Our simulator of the Mueller--D'Andrea drone reproduces their two-rotor and three-rotor-loss equilibria to every published digit, five hover solutions of their more detailed model within 4\%, and the spin-up time of their flight within about 4\%. For prediction horizons from 100\,ms to 1\,s, the model that learns only the error (the \emph{residual} model) is the most accurate predictor on the plant it was trained on, and with 10 training episodes it beats a full Koopman model trained on 40. Inside model predictive control (MPC) it reduces a steady altitude error caused by model mismatch by about one third. The learned models do not yet work on a plant with different parameters. In our simulation, the published recovery strategy recovers safely with probability above 0.9 only if the failure is detected within about 12\,ms, and below 0.5 beyond about 19\,ms, while published detectors need about 20 to 130\,ms. A first certificate, valid for the linearized model with thrust limits and based on linear matrix inequalities, covers 22\% of the way from the spinning equilibrium to the state right after a failure.

Technical terms are explained where they first appear and are collected in a glossary (Appendix~\ref{app:glossary}). Each method is described in three layers: the idea in plain words, the mathematics, and how it is implemented in our code.
\end{abstract}
